\scanpage{28}
On a donc une \uline{extension}
\diagram{\begin{tikzcd}[column sep=small]
1\arrow[r]&G'_S(\unsure{S|\Omega})\arrow[r,dash]&\mathcal E\arrow[r]&\Gamma\arrow[r]&1\arrow[r,dash]&\unsure{0|\circ}
\end{tikzcd}}
\hfill $(\mathcal E)$

\deleted{Supposons que $\pi\unsure{\in|\subset}\Gamma$ opère sur}
\inserted{$\pi$}

Les \uline{sections} de celle-ci sont les structures de torseur sur $(S,\Gamma;G_S)$ sur $Q_S$ qui
\unsure{prolongent} sa structure de tor. sur $(S,G_S)$.
Si on a une telle section, on en déduit une opération de $\Gamma$ sur $G'_S(\unsure{S|\Omega})$
(qui en fait provient de son opération sur $G'_S$ lui-$=$ !)
et les autres sections correspondent aux 1-cocycles de $\Gamma$ à valeurs dans $G'_S(\unsure{S|\Omega})$.
\deleted{Les}\inserted{L'ens $\Gamma(\mathcal E)$ des} sections
forment \deleted{\illegible{1}} un ens où $G'_S(\unsure{S|\Omega})$ opère
(\unsure{autom.|automs} intérieurs) et deux \deleted{\unsure{de pts. de|élts. de}}
$\Gamma(\mathcal E)$ sont
\inserted{conjugués ssi le 1-cocycle faisant passer de l'un à l'autre
\unsure{est coh. à} $1$}.
\figure{\begin{tikzpicture}[x=1cm,y=1cm]
\draw (0,0) .. controls (-0.15,0.4) and (1.7,0.7) .. (2,0.8);
\draw (0.1,0.5)--(0.3,0.3)--(0.4,0.55);
\end{tikzpicture}}

Si on a un \unsure{ss-gr|sous-gr} $\pi\subset\Gamma$ et une section
$\alpha_0:\pi\longrightarrow\mathcal E$ partielle, on peut s'intéresser
à l'ens des $g\in G'_S(\unsure{S|\Omega})$ qui la fixent, c'est \unsure{l'ens} des
élts de $G'_S(\unsure{S|\Omega})$ fixés par $\pi$ (y opérant comme de juste)
$G'_S(\unsure{S|\Omega})^\pi\ (=G'_L(L))$.
Les sections $\alpha$ qui prolongent $\alpha_0$ sont donc telles que
$G'_S(\unsure{S|\Omega})^\pi$ y opère, Si on a
\margin{27}
